
\definecolor{bdcolor}{RGB}{68,170,168} % define border color
\newcommand{\overviewBox}[1]{%
\tikz[baseline=-0.5ex]{\node[rounded corners=0.2em, 
fill=white, 
draw=bdcolor, 
minimum size=0.9em, 
text=bdcolor, 
text centered,
inner sep=0pt,
line width=0.9pt,
font=\fontsize{9pt}{0.9em}\fontseries{ul}\selectfont](TsNode){#1}}}


\section{Introduction}
Large Language Model (LLM) based agents, which are capable of observing, making decisions, and performing actions with the reasoning capabilities of LLM, have undergone significant improvements, showing great potential in various areas such as programming\cite{hong2023metagpt,chetDev}, creative writing\cite{AutoAgents, wang2023unleashing, OKRagents}, and question answering \cite{wu2023empiricalMath, tang2023MedAgents}. 
% Inspired by the synergistic effects brought about by human cooperation
While initial forays into the realm of LLM-based agents focused on solitary agent system\cite{AutoGPT,Expertprompting}, the concept of multi-agent collaboration—mirroring the cooperative interactions among humans—has started to pique the interest of the AI research community. Drawing inspiration from the synergistic outcomes observed in human teamwork\cite{engelbart2023augmenting, woolley2010evidence,luppi2022synergistic}, a burgeoning body of research has been investigating and validating the benefits (e.g. expand expertise\cite{Salewski2023InContextIR, tang2023MedAgents}, enhance reliability \cite{chan2023chateval, du2023improving_debate, chetDev}, encourage divergent thinking\cite{Encouraging_Divergent_Thinking_Debate, zhuge2023mindstorms}) brought about by LLM-based multi-agent collaboration.

In order to facilitate the coordination of multi-agent collaboration, the open-source community emerges a multitude of frameworks for prototyping LLM-based multi-agent systems. Existing multi-agent frameworks can be divided into two categories based on how users can specify or intervene in the collaborative process: code-based and natural language-based. For code-based frameworks \cite{hong2023metagpt, chetDev, wu2023autogen, CrewAI, Langroid, chen2023agentverse}, users need to hard-code the coordination strategies (e.g. the division of tasks, the assignment of agents, the flow of massages) directly into the code, which requires code skill and learning cost. The natural language-based frameworks\cite{AutoAgents, wu2023autogen} \footnote[1]{{AutoGen \cite{wu2023autogen} supports both code-based and natural language-based paradiam. In its ``group chat mode'', the coordination strategy can be expressed in free-form natural language and coordinated by a chat manager.}}, which directly uses natural languages to specify the coordination strategies, could be a promising way to democratize agent technology for broader general users. Additionally, given that LLMs inherently possess coordinating capabilities and rich domain knowledge across different tasks, the natural language-based frameworks can easily leverage LLMs to assist in drafting and refining the coordination strategies represented in natural language\cite{AutoAgents}. 

However, designing coordination strategies remains challenging with existing natural language-based frameworks. First, the flexibility of natural language could be a double-edged sword: on one hand, it allows users to freely design and express their coordination strategies; on the other hand, overly flexible expressions can make the devised collaboration strategies ambiguous, often requiring users repeatedly engage in remedial specification to ensure the execution of the collaboration doesn't stray from its intended course. Second, representing and exploring coordination strategies in pure text format faces challenges as the complexity of the collaboration process and team organization rises. Users can easily get lost in the ``text jam'' where important information (e.g. agent relationship, task dependency, result correspondence, strategy discrepancy) they care about at certain points during exploration is drowned in blocks of text. Therefore, novel approaches are highly desirable to enhance the current natural language-based design process.

This work thus presents a visual exploration framework for efficiently designing coordination strategies for LLM-based multi-agent collaboration. To exploit the flexibility of natural language while incorporating a level of organization to regularize its ambiguity, we analyze the common concepts and structures found in the description for coordination strategy in a corpus of 25 LLM-based multi-agent collaboration papers and 7 high-star projects \footnote[2]{{The corpus can be found in our project repository.}}, based on which we establish a structured representation for LLM-based Multi-agent coordination strategy. This structure lays a foundational scaffolding for the entire exploration process. Based on this structure, we design a generation method that exploits the coordinating capabilities and domain knowledge of LLMs to map the general goal provided by users (\cref{fig:teaser} \overviewBox{a}) into an executable initial coordination strategy to help users kick off the exploration. To ensure coherence among various parts of the generated strategy, we divide the generation process into three stages: Plan Outline Generation (\cref{fig:teaser} \overviewBox{b}) for an overall collaboration plan to achieve the goal, Agent Assignment (\cref{fig:teaser} \overviewBox{c}) for each task in the plan outline, and Task Process Generation (\cref{fig:teaser} \overviewBox{d}) for specifying how assigned agents collaboratively finish the task. To streamline users' exploration and iterative refinement for alternative strategies, we propose a set of interactions (\cref{fig:teaser} \overviewBox{e} \overviewBox{f} \overviewBox{g}) to help users visually explore the design space for each generation stage with the help of LLMs. Whenever users are satisfied with a certain strategy, they can initiate the collaboration and examine the execution result which is visually enhanced and linked with previous stages for efficient verification (\cref{fig:teaser} \textbf{\overviewBox{h}}). 

To validate the feasibility and effectiveness of this framework, we developed an interactive system called \textit{AgentCoord} that enables users to visually explore coordination strategies for LLM-based multi-agent collaboration, effectively integrating the prior of LLMs and users during design. Our user study, involving 12 users with a general interest in LLM-based multi-agent collaboration, suggests that our approach can effectively facilitate the design process for LLM-based multi-agent coordination strategies and has the potential to democratize agent coordination for broader users. In summary, our contributions include:

\begin{itemize}
\item A visual exploration framework that enables general users to efficiently design
coordination strategy for LLM-based multi-agent collaboration.
\item \textit{AgentCoord} \footnote[3]{{Project Repository: https://github.com/AgentCoord/AgentCoord}}, an open-source interactive system that instantiates our framework with a set of interactions and visual designs to facilitate coordination strategy exploration.
\item A formal user study that
demonstrates the feasibility and effectiveness of our approach.
\end{itemize}







